% --- Back cover ---
% Large PERSONIX wordmark, no logo image. Cream paper inherited globally.
% Absolute (tikz overlay) positioning so the footer sticks to the bottom.
\clearpage
\thispagestyle{empty}
\begin{tikzpicture}[remember picture, overlay]
  % --- Wordmark ---
  \node[anchor=north, font=\sffamily\bfseries\fontsize{48}{52}\selectfont]
    at ([yshift=-26mm]current page.north) {PERSONIX};
  \node[anchor=north, font=\rmfamily\itshape\large]
    at ([yshift=-44mm]current page.north) {Бескомпромисна промена};
  \node[anchor=north, font=\rmfamily\small,
        text width=\paperwidth-50mm, align=center]
    at ([yshift=-54mm]current page.north)
    {Нецензурабилна и непоткуплива децентрализирана репутациска мрежа};
  % --- Body ---
  \node[anchor=north, text width=\paperwidth-50mm, align=justify, font=\small]
    at ([yshift=-72mm]current page.north)
    {%
      Државата функционираше кога комуникацијата беше бавна, одлуките беа
      локални, а власта живееше во истиот град со кој владееше. Денешните
      мрежи ги инвертираат сите три претпоставки --- и институциите изградени
      врз нив веќе не му одговараат на светот со кој управуваат. Оваа книга
      опишува алатка што одговара: децентрализирана репутациска мрежа во која
      идентитетот е твој да го задржиш, вистината е јавно проверлива, а
      поткупот е репутациско самоубиство.

      \smallskip
      Не манифест. Не прогноза. Работен нацрт за тоа како може да се изгради
      наследникот на државата --- доброволно, една декларација наеднаш.
    };
  % --- Footer (anchored to the bottom of the page) ---
  \node[anchor=south, font=\sffamily\footnotesize,
        text width=\paperwidth-50mm, align=center]
    at ([yshift=12mm]current page.south)
    {%
      Pavel Kudrna\quad\textbullet\quad Прага, 2026\par
      \href{https://personix.org}{personix.org}\par
      Лиценцирано под Creative Commons Attribution-ShareAlike 4.0 International
      (CC BY-SA 4.0).
    };
\end{tikzpicture}%
\mbox{}%
\clearpage
